\author{OTS}
\documentclass[fontsize=16pt]{mybib}
\setheaderlogo{logo.png}

% Absatzformatierung
\setlength{\parindent}{0pt}
\setlength{\parskip}{0.6em}
\renewcommand{\arraystretch}{1.8}
\renewcommand{\DFA}{12mm}
\renewcommand{\DZA}{1.2}

\title{\biblezit{IIPetr}{3:1-8}}
\author{Lothar Schmid}
\date{}

\begin{document}
\mytitlepage
\tableofcontents

\section{2. Petrus 3}

\section{Einleitung}
\begin{center}
  \begin{verse}
  Alle Jahre wieder\\
  Kommt das Christuskind\\
  Auf die Erde nieder,\\
  Wo wir Menschen sind;    
  \end{verse}
\end{center}
Ein bekanntes Kinderlied komponiert von Friedrich Silcher. Ich denke viele hier haben jetzt sicher die Melodie Kopf. Ich werde das jetzt sicherheitshalber nicht vorsingen.

Alle Jahre wieder, am 25 Dezember feiern wir Weihnachten. Der 25. Dezember ist kein Zufall, sondern wurde von Hippolyt von Rom festgelegt. Er schreibt: 
\begin{quote}\enquote{Die erste Ankunft unseres Herrn ... geschah in Bethlehem am achten Tag vor den Kalenden (Kalendae wurde der erste Tag im Römischen Kalender genannt, daraus entstand das Deutsche Wort Kalender).} 
\end{quote}
Nach römischer Zählung entspricht das dem 25. Dezember. Das Datum wurde aber nicht einfach PI mal Daumen hingeklatscht, sondern das wurde berechnet. Nach der alten jüdischen Tradition stirbt ein grosser gerechter Mensch an dem Tag, an dem er gezeugt wurde. Also wenn Jesus am 14. Nisan oder 25. März, am Kreuz gestorben ist, dann war seine Zeugung auch an diesem Tag und so neun Monate später seine Geburt.

Lange Zeit hatte aber der Geburtstag Jesu keinen grossen Stellenwert. Er wurde auch nicht speziell gefeiert. Johannes Chrysostomus (ca. 349-407) erwähnt, dass die Feier am 25. Dezember in Antiochia noch relativ neu sei, verteidigt sie aber als römische Überlieferung.

In der ganzen Bibel lesen wir nirgends, dass die Juden Geburtstage gefeiert haben. Weder von den alten Erzvätern noch von David oder Salomo. Es gibt nur zwei Erwähnungen von Geurtstagsfeiern in der Bibel und beide male von heidnischen Königen. Einmal den Geburtstag des Pharao in \biblezit{IMos}{40:20} und das zweite mal in \biblezit{Matt}{14:6} oder \biblezit{Mk}{6:21}. Die bekannte Geschichte, bei der Johannes der Täufer sterben musste.

Die Adventszeit entsteht aber viel später. Die ersten Formen einer Vorbereitungszeit treten im 5. Jahrhundert in Spanien auf. Es war eine Art fastenzeit wie vor Ostern.

Gregor der Grosse (ca. 590-604) führt dann die römische Form von vier Adventssonntagen ein. Spannend ist, dass der frühchristliche Advent stark auf das zweite Kommen Jesu ausgerichtet war und weniger auf das erste Kommen Jesu. So wurde an diesen Adventssonntagen oft aus Matthäus 24, Markus 13 und Lukas 21 vorgelesen, wo vom zweite Kommen Jesu die Rede ist. 

Mit der Zeit verschiebt sich der Fokus immer mehr auf den Geburtstag von Jesus. Die Gebräuche, die wir heute haben, sind alle ziemlich neu. Der erste Adventskranz, so um 1870, wurde von Johann Hinrich Wichern aufgehängt. Dieser Kranz war ein an der Decke aufgehängtes Wagenrad mit 25 Kerzen. Es hing im sogenannten \enquote{Rauhen Haus} in der Nähe von Hamburg, einem Kinderheim für arme Kinder. Es sollte diesen armen Kinder die Vorfreude auf Weihnachten erhöhen. Aus diesem Wagenrad wurde dann ein Kranz mit vier Kerzen, den man einfach auf den Tisch stellen konnte.

Der erste gedruckte Adventskalender kam 1902 raus auch in Hamburg.

Ich habe noch zwei Schwestern und der, der den ganzen Tag brav war, durfte am Abend das Türchen vom Adventskranz aufmachen. Am Abend haben wir dann darum gestritten, wer denn nur der Bravste tagsüber war und das Türchen aufmachen darf. Ich sage jetzt nicht, wer meistens gewonnen hat.

Inzwischen ist es schwierig ziemlich geworden, in einem Geschäft einen Adventskalender mit christlichen Motiven zu bekommen. Am besten mit den Kindern oder Enkel zusammen selber einen basteln.

\section{Die Erinnerung}
Nun aber! Worauf warten wir heute im Advent? Um darauf eine Antwort zu bekommen ist das dritte Kapitel im zweiten Petrus Brief sehr gut geeignet. Der Brief von Petrus spricht in unsere Zeit rein. Er ist unglaublich aktuell. Lesen wir die ersten beiden Verse.
\begin{bibelbox}{SCHL}{IIPetr}{3:1-2}
  Geliebte, dies ist nun schon der zweite Brief, den ich euch schreibe, um durch Erinnerung eure lautere Gesinnung aufzuwecken, damit ihr die Worte gedenkt, die von den heiligen Propheten vorausgesagt worden sind und dessen, was euch der Herr und Retter durch uns, die Apostel aufgetragen hat.
\end{bibelbox}
Wie wir hier lesen, ist dies der zweite Brief den Petrus an seine Geschwister schreibt. Aus dem ersten Brief wissen wir, dass Petrus seine Briefe an jüdische Christen in der Diaspora, oder Zerstreuung, im heutigen Süden der Türkei gerichtet hat. Im ersten Brief betont er mehr die Ermutigung und wie sie sich in der heidnischen Umgebung verhalten sollen. Hier in diesem zweiten Brief ermahnt er sie, sich nicht beeinflussen und vom Glauben abbringen zu lassen.\\
Und so wie ich das verstehe, nicht nur von den Heiden, sondern vor allem von solchen, die vorgeben die Schriften zu kennen. In Vers 4 schreibt Petrus:
\begin{bibelbox}{SCHL}{IIPetr}{3:4}
  Wo ist die Verheissung seiner Wiederkunft?
\end{bibelbox}
und in den weiteren Versen fünf uns sechs nimmt Petrus Bezug auf die Schöpfung und auf die Sintflut und geht davon aus, dass diese Spötter die Geschichten aus der Bibel kennen \enquote{Dabei übersehen sie aber absichtlich ...} wie hier in Vers 5 angeführt.

Es ist für uns sehr wichtig, dass auch wir uns in der Schrift gut auskennen. Wenn wir heute auf den sozialen Medien unterwegs sind und Predigten suchen, ist die Gefahr gross, dass wir auf die Falschen treffen. Als ich noch jung im Glauben war, fand ich den Kanal \enquote{Worthaus}. Aktuell hat dieser Kanal ca. 15600 Abonnenten und 337 Vorträge.\\
Die Beschreibung liest sich so Zitat:
\begin{quote}
  Der Worthaus e. V. wurde 2010 mit dem Ziel gegründet den aktuellen Diskussionsstand der christlichen Hochschultheologie einem breiten Publikum verständlich und in Form von Video- und Audiovorträgen im Internet kostenfrei zugänglich zu machen. Worthaus will allen, die sich für den christlichen Glauben interessieren, ein spannendes, informatives und inspirierendes Angebot auf hohem inhaltlichem Niveau machen. Dabei legt Worthaus Wert auf theologische Substanz, Verständlichkeit und gesellschaftliche Relevanz.
\end{quote}
Ich habe zwei Vorträge von Prof. Dr. Siegfried Zimmer gehört. Ich fand diese so gut, dass ich sie sie sogar meiner Frau vorgeschlagen habe. Bei dem 3. Vortrag wurde ich skeptisch. Er redete von Allversöhnung, Gott ist kein strafender Gott usw. Da wurde ich stutzig und habe mich dann ein kritischer mit den Inhalten auf diesem Kanal befasst und unter anderem ein Video gefunden, in dem er die Christen lächerlich macht, die Jesus als Gottes Sohn ansehen, an die Auferstehung und an die Jungfrauengeburt glauben.

Ein Spötter, wie Petrus ihn beschreibt. Einer, der die Bibel kennt und den Inhalt ins Lächerliche zieht. Wir können das nur erkennen, wenn wir das wahre Wort Gottes kennen, es studiert haben und auch Glauben, dass dieses Buch hier von Gott inspiriert ist.

Jesus sagt in \biblezit{Mk}{13:5} nicht nur, dass Spötter auftreten werden, sondern dass gegen Ende immer mehr von ihnen auftreten werden. Er war mit Petrus, Jakobus, Johannes und Andreas zusammen, als er ihnen sagte: \enquote{Habt acht, dass euch niemand verführt!} Diese Mahnung will Petrus seinen Geschwistern in der Diaspora weitergeben.

In Vers 4 bringt Petrus eine so typsche Aussage von Skepitkern oder Spöttern.
\begin{bibelbox}{SCHL}{IIPetr}{3:4}
\enquote{Wo ist die Verheissung seiner Wiederkunft? Denn seitdem die Väter entschlafen sind, bleibt alles so wie es von Anfang der Schöpfung gewesen ist.}
\end{bibelbox}
Habt ihr diese Frage nicht auch schon gehört? Normalerweise wollen die Fragesteller keine Antwort darauf hören. Sie stellen sie als eine sogenannte rethorische Frage. Also ein Frage auf die keine Antwort erwartet wird. Das erste Kommen von Jesus haben sie verpasst. Eigentlich sollte ja der kommende Messias ein neues Königreich aufbauen. Aber nichts hat sich geändert. Die Römer sind noch da, in Jerusalem brodelds es immer mehr, und diese komischen Christen reden da irgendwas von einem Erlöser, der wiederkommen wird.

Wer heute vom zweiten Kommen Jesu spricht, wird auch ausgelacht, oder als Endzeitfanatiker hingestellt. Im Frühjahr kam im deutschen Fernsehen eine Dokumentation über junge Christen. Es hiess dort: \enquote{Hier handelt es sich um eine rechtsradikale, fundamentalistische Gruppe, die auf die Endzeit wartet.}\\
Es gab in der ganzen Sendung auch nur zwei Interview, in denen die Jugentlichen erklären, dass sie glauben, dass Jesus am Ende der Zeit wieder kommt und Gericht hält.

Diese Ungeduld, wann wirkt Gott endlich, ist nicht neu. So lesen wir in \bibleverse{Hes}(12:22-23) folgendes.
\begin{bibelbox}{SCHL}{Hes}{12:22-23}
\hh{22}Menschensohn, was ist das für ein Sprichwort, das ihr im Land Israel gebraucht, indem ihr sagt: \enquote{Die Tage ziehen sich hinaus, und es wird nichts aus allen Offenbarungen}\\
\textit{Darauf antwortet Gott im Vers 23b}\\
\hh{23} Du aber sprich zu ihnen: \enquote{Die Tage sind nahe, und jedes Wort der Offenbarung [trifft bald ein]!}
\end{bibelbox}

Für die Menschen, die von Petrus angeschrieben wurden, war das Christenleben hart. Verfolgung, Spott und Unterdrückung war ihr Alltag. Sie sehnten sich nach Frieden und Freiheit, die sie nur bekommen würden, wenn Jesus wieder kommt. Mit der Wiederkunft Jesu haben sie sich gegenseitig getröstet und sich Mut zugesprochen.

Hier im Mitternachtsruf werden die Gebete oft mit dem Ausdruck \betonung[b, p]{Maranatha} beendet, was soviel heisst wie \enquote{komme bald}. Wünschen wir uns das wirklich von Herzen? Oder doch lieber noch bis nach Weihnachten warten? Aktuell habe ich noch keine Zeit. Ich will doch wissen, was ich zu Weihnachten bekomme.

Ich will damit sagen, dass es uns hier im Westen noch sehr gut geht und das finde ich auch super so und wir sollen unserem Herrn auch dankbar dafür sein. Wir haben es uns ja nicht ausgewählt in welchem Land wir diese Erde betreten. Ist hier zu leben, nicht ein Grund unserem Gott dauernd zu danken?\\

Die Christen vor 1900 Jahren warteten auf ihre Erlösung aus ihrem Elend und jetzt kommen da Leute, wir würden sie heute Theologen nennen, die ihnen erzählen, dass sie vergeblich warten. Oft sind es noch Menschen die auch ohne Glauben erfolgreich sind. Wie auch Asaf in seinem Psalm 73 ausruft:
\begin{bibelbox}{SCHL}{Ps}{73:3}
  Denn ich beneidete die Übermütigen, als ich das Wohlergenen der Gottlosen sah.
\end{bibelbox}
Das kennen wir doch auch, oder?\\
Darum weist Petrus in seinem zweiten Brief auch wieder auf die Sindflut hin. Petrus hat dieses Beispiel schon ein paar Zeilen vorher in \biblezit{IIPetr}{2:5} und auch in seinem ersten Brief \biblezit{IPetr}{3:20} beutzt wo er schreibt:
\begin{bibelbox}{SCHL}{IPetr}{3:20}
die vor Zeiten sich weigerten zu glauben, als Gottes langmut erst mal zuwartete in den Tagen Noahs, während die Arche zugerichtet wurde, in der wenige, nähmlich acht Seelen, hindurchgerettet wurden.
\end{bibelbox}
Die Spötter wollen nicht wahrhaben oder wollen vergessen, dass Gott handelt. Die Theologen und Bibelkenner wissen, was in der Bibel steht. Aber erklären es als Geschichte als Märchen und Aberglaube. Lasst euch aber von denen nicht verführen. Der Widersacher gebraucht alles, um uns von Jesus wegzuziehen. Vorallem um zu verhindern, dass andere Leute die Jesus nicht kennen, ihn kennen lernen. Wenn du Jesus kennst und ihn angenommen hast, dann bist du versiegelt und Satan kann dich nicht mehr von ihm wegreissen. Aber du kannst den Geist betrüben und dich so immer weiter von Jesus entfernen.\\
\section{Der Trost}
In Vers 7 betont Petrus nochmals das Versprechen, das Gott Noah nach der Sintflut geben hat. 
\begin{bibelbox}{SCHL}{IMos}{9:11}
Und ich will meinen Bund mit euch und mit eurem Samen, der nach euch kommt, auch mit allen lebendigen Wesen bei euch, aufrichten, dass künftig nie mehr alles Fleisch von dem Wasser der Sintflut ausgerottet wird, und dass keine Sintflut mehr kommen soll, um die Erde zu verderben.
\end{bibelbox}
Das nächste Gericht, das Gott ankündigt, wird also keine Flut mehr sein, sondern Feuer.
\begin{bibelbox}{SCHL}{IIPetr}{3:7}
Die jetzigen Himmel aber und die Erde werden durch dasselbe Wort aufgespart und für das Feuer bewahrt bis zum Tag des Gerichts und des Verderbens der gottlosen Menschen.
\end{bibelbox}
Vers 8 ist einer meiner Lieblingsverse, weil man ihn so einfach auswendiglernen kann. Aber anscheindend nicht nur von mir, sondern auch von denen, die nie die Bibel gelesen haben.
\begin{bibelbox}{SCHL}{IIPetr}{3:8}
Dieses eine aber sollt ihr nicht übersehen, Geliebte, dass ein Tag bei dem Herrn ist wie tausend Jahre, und tausend Jahre wie ein Tag.
\end{bibelbox}
Mit diesem Vers wird oft das Âlter der Erde in der Schöpfungsgeschichte erklärt, und das sogar auch von vielen Christen. Wenn ihr diesen Vers nehmt und dafür gebraucht um aus 6 Schöpfungtagen drei Milliarden Jahre zu machen, dann bitte ich euch, diesen Vers mal umzudrehen. Nicht tausend Jahre wie eine Tag, sondern 1 Tag wie tausend Jahre.

Gott sprach und der Baum stand, also nicht mal eine Sekunde. Hätte man ihn danach direkt gefällt, so hätte man an seinen Jahresringen zählen können, dass er schon über 100 Jahre alt war. Wie bei der Hochzeit von Kana. Die Gäste mussten auch  nicht ein Jahr warten bis der Wein greift war.

Im Vers 8 geht es also nicht um eine Zeitangabe, sondern darum, dass Gott nicht unserer irdischer Zeitrechnung unterliegt. Petrus zitiert hier aus \biblezit{Ps}{90:4}
\begin{bibelbox}{SCHL}{Ps}{90:4}
Denn tausend Jahre sind vor dir wie der gestrige Tag, der vergangen ist,
\end{bibelbox}

Natürlich ist das ein schwacher Trost, wenn es einem schlecht geht. Darum gibt Petrus in Vers 9 eine Erklärung dazu ab
\begin{bibelbox}{SCHL}{IIPetr}{3:9}
Der Herr zögert nicht die Verheissung hinaus, wie etliche es für ein Hinauszögern halten, sondern er ist langmütig gegen uns, weil er nicht will, dass jemand verlorengehe, sondern dass jedermann Raum zur Busse bekomme.
\end{bibelbox}

Jeder von uns hat doch sicher Angehörige, die Jesus noch nicht kennen oder kennen wollen. Wir müssen uns bewusst sein, dass wenn jetzt in diesem Augenblich Jesus wiederkommt, alle in der Hölle landen. Ich glaube aber an die Entrückung und hoffe, dass sie wenigsten nach diesem Ereignis zur Besinnung kommen und sagen: \enquote{War doch kein Unsinn was Lothar uns da erzählt hat.}

Verse 10-13
\begin{bibelbox}{SCHL}{IIPetr}{3:10-13}
\hh{10}Es wird aber der Tag des Herrn kommen wie ein Dieb in der Nacht; dann werden die Himmel mit Krachen vergehen, die Elemente aber vor Hitze sich auflösen und die Erde und die Werke darauf verbrennen.\\
\hh{11}Da nun dies alles aufgelöst wird, wie sehr solltet ihr euch auszeichnen durch heiligen Wandel und Gottesfurcht,\\
\hh{12}indem ihr das Kommen des Tages Gottes erwartet und ihm entgegeneilt, an welchem die Himmel sich in Glut auflösen und die Elemante vor Hitze zerschmelzen werden!\\
\hh{13}Wir erwarten aber nach seiner Verheissung neue Himmel und eine neue Erde in denen Gerechtigkeit wohnt.
\end{bibelbox}
Das lesen wir auch in der Offenbarung \biblezit{Offb}{21:1}.
\begin{bibelbox}{SCHL}{Offb}{21:1}
  Und ich sah einen neuen Himmel und eine neue Erde; denn der erste Himmel und die erste Erde waren vergangen, und das Meer gibt es nicht mehr.
\end{bibelbox}
Intressant ist, dass Petrus als er diesen Brief schrieb noch keinen Schimmer hatte, das Johannes irgendwann diese Visionen von Jesus bekommt. Hier ist nochmals ein Beweis, dass alles was in diesem Buch steht nicht Menschenwort ist, sondern Gottes Wort.\\
Wir wissen nicht wann Jesus wiederkommt. Aus diesem Grund sollen wir stehts vorbereitet sein. Darum schreibt Petrus in Vers 1
\begin{bibelbox}{SCHL}{IIPetr}{3:1}
und durch Erinnerung eure lautere Gesinnung aufzuwecken.
\end{bibelbox}
Und genau darum ist dieser Brief auch in der Bibel. Petrus hat ihn damals seinen Glaubensgeschwister geschrieben. Gott hat ihn in den Kanon der Schrift aufnehmen lassen und uns aufzuwecken. Es ist nicht eine Prüfung Gottes oder er will uns hier nicht einfach so hinhalten. Uns geht es gut, aber für die verfolgten Christen?\\

\section{Das Erwarten}
Also auf was warten wir jetzt im Advent? Vor 1900 Jahren habe die Christen auf das Wiederkommen Jesus erwartet. Wir wissen, dass Jesus noch nicht kommt, weil die Vollzahl der Heiden noch nicht erfüllt ist und weil die Liebe und Gnade Gottes nicht will jemand verloren geht.
Wenn du Jesus jederzeit erwartest, hast du keine andere Wahl als so zu Leben, dass wenn Jesus wiederkommt die als Tadellos vorfindet.
\begin{bibelbox}{SCHL}{IIPetr}{3:14}
Darum, Geliebte, weil ihr das ihr dies erwartet, so seid eifrig darum bemüht, dass ihr als unbefleckt und tadellos vor ihm erfunden werdet in Frieden!
\end{bibelbox}

\begin{bibelbox}{SCHL}{IIPetr}{3:15-16}
Und seht die Langmut unseres Herrn als [eure] Rettung an, wie auch unser geliebter Bruder Paulus euch geschrieben hat nach der ihm gegebenen Weisheit, so wie auch in allen Briefen, wo er von diesen Dingen spricht. In ihnen ist manches schwer zu verstehen, was die Unwissenden und Ungefestigten verdrehen, wie auch die übrigen Schriften ihrem eigenen Verderben.
\end{bibelbox}
\begin{bibelbox}{SCHL}{IIPetr}{3:17}
Ihr aber Geliebte, da ihr dies im Voraus wisst, so hütet euch, dass ihr nicht durch die Verführung der Frevler mit fortgerissen werdet und euren eigenen festen Stand verliert!
\end{bibelbox}
\begin{bibelbox}{SCHL}{IIPetr}{3:18}
Wachst dagegen in der Gnade und in der Erkenntnis unseres Herrn und Retters Jesus Christus! Ihm sei die Ehre, sowohl jetzt als auch bis zum Tag der Ewigkeit! Amen
\end{bibelbox}
\betonung[b, p]{Erlebniss Adventsmarkt}



























\end{document}